\begin{table}[htbp]
\centering
\caption{Standardized Effect Sizes}
\label{tab:sde}
\begin{minipage}{\textwidth}
\centering
\begin{adjustbox}{max width=\textwidth}
\begin{tabular}{lcccccc}
\toprule
Outcome & $\hat{\beta}$ & SE & SD($Y$) & SDE & SE(SDE) & Classification \\
\midrule
\multicolumn{7}{l}{\textit{Panel A: Pooled}} \\
Female non-mining emp (boom) & -0.1220 & 0.0534 & 1.5471 & -0.0789 & 0.0345 & Moderate negative \\
Female non-mining emp (bust) & -0.1171 & 0.0641 & 1.5471 & -0.0757 & 0.0414 & Moderate negative \\
Female non-mining earn (boom) & -0.0767 & 0.0187 & 0.2271 & -0.3375 & 0.0822 & Large negative \\
Gender earnings gap (boom) & 0.0666 & 0.0130 & 0.1158 & 0.5751 & 0.1122 & Large positive \\
\midrule
\multicolumn{7}{l}{\textit{Panel B: Heterogeneous (Urban vs.\ Rural)}} \\
Female non-mining emp (urban) & -0.0447 & 0.0418 & 1.0860 & -0.0412 & 0.0385 & Small negative \\
Female non-mining emp (rural) & -0.2159 & 0.0673 & 0.9366 & -0.2305 & 0.0719 & Large negative \\
\bottomrule
\end{tabular}
\end{adjustbox}

\vspace{0.5em}
\footnotesize
\textit{Notes:} \textbf{Country:} United States. \textbf{Research question:} Does male-biased labor demand from the shale oil and gas boom reduce female non-mining employment and earnings in affected counties? \textbf{Policy mechanism:} The shale fracking revolution created massive mining employment that is approximately 14:1 male-to-female, generating a male-biased labor demand shock that distorted local labor markets through in-migration, cost-of-living increases, and industry composition shifts in affected counties. \textbf{Outcome definition:} Log average quarterly non-mining employment count for female workers from the Quarterly Workforce Indicators (QWI), and gender earnings gap computed as (male minus female non-mining earnings) divided by male earnings. \textbf{Treatment:} Binary indicator for counties in the top quartile of pre-boom (2001--2005) mining employment share. \textbf{Data:} QWI county-year-sex-industry panels from Census Bureau, 24 states, 2001--2022, 59,090 county-year-sex observations covering 249 treated counties. \textbf{Method:} Triple-difference (female $\times$ high-mining $\times$ boom period) with county and year fixed effects; standard errors clustered at state level. \textbf{Sample:} Counties in 24 states spanning major shale plays (Bakken, Permian, Eagle Ford, Marcellus) and comparison states; restricted to counties with non-missing QWI data. SDE $= \hat{\beta} / \text{SD}(Y)$ where SD($Y$) is the pre-treatment (2001--2005) standard deviation. Classification refers to magnitude, not statistical significance: Large ($|$SDE$| > 0.15$), Moderate ($0.05$--$0.15$), Small ($0.005$--$0.05$), Null ($< 0.005$).
\end{minipage}
\end{table}
